\section{Genuine tensor multiplicativity of original proper-model line norms}

\begin{definition}[Actual tensor frame on the original proper center]
\label{def:actual-center-line-tensor-frame}
\lean{AlgebraicGeometry.Scheme.Modules.IsInvertible.monoidalTensor}
\lean{EnvelopeProperty.Berkovich.RelativeSchemeValuedPoint.centerTensorFrame}
\leanok
\uses{lem:scheme-line-actual-center-frame}
Given actual original center frames of two original invertible
modules, use the inverse native monoidal tensorator, the tensor of
these actual frames, and the actual structure-module multiplication
isomorphism to frame their actual pulled tensor line. The explicit
sheafified tensor and the native monoidal tensor are identified by
the genuine original comparison isomorphism.
\end{definition}

\begin{lemma}[Actual generic tensor frame]
\label{lem:actual-center-generic-tensor-frame}
\lean{EnvelopeProperty.Berkovich.RelativeSchemeValuedPoint.centerFrameGenericIso_tensor}
\leanok
\uses{def:actual-center-line-tensor-frame,lem:scheme-line-actual-generic-frame-naturality}
The actual generic frame of this original center tensor frame is
the inverse original generic tensorator followed by the tensor of
the two original generic frames and the original structure-module
multiplication isomorphism.
\end{lemma}
\begin{proof}
The genuine composition comparison of actual scheme pullbacks is
monoidal, and so is the genuine comparison for the actual proper
valuative generic triangle. The original actual tensor-frame
composition theorem applies to their composite. Its two pulled
factor frames are exactly the original generic frames by their
actual definitions; its other side is exactly the original generic
frame of the actual center tensor frame. This gives the asserted
whole original frame equality without a coherence premise.
\end{proof}

\begin{theorem}[Original pure tensor section norm]
\label{thm:actual-model-line-tensor-norm}
\lean{EnvelopeProperty.Berkovich.RelativeSchemeValuedPoint.modelLineNorm_tensor}
\lean{EnvelopeProperty.Berkovich.RelativeSchemeValuedPoint.genericTensorSection}
\lean{EnvelopeProperty.Berkovich.RelativeSchemeValuedPoint.modelLineNorm_genericTensorSection}
\leanok
\uses{lem:actual-center-generic-tensor-frame,def:scheme-line-actual-model-norm,lem:scheme-line-actual-model-isomorphism-isometry}
The original proper-model norm of the actual pure tensor of two
sections on the same original residue-field line fibers is the
product of their original model norms, for both the native
monoidal tensor and the actual explicit sheafified tensor.
\end{theorem}
\begin{proof}
Choose original center frames supplied by the actual original
invertibility. Compute the original tensor-line norm in the actual
center tensor frame. By the preceding whole frame equality, its
coefficient is obtained by applying the two original generic
frames to the factors and multiplying the two resulting original
structure sections: naturality of the true original pure tensor
section and actual structure-module multiplication give this exact
coefficient identity. The actual global-section isomorphism and
the original residue valuation both preserve multiplication. Thus
the coefficient value is the product of the original factor norms.
Finally the original coherent-to-explicit tensor comparison is an
actual whole original line-module isomorphism, hence an isometry
for the genuine original model norms. It carries the constructed
native pure tensor to the stated actual explicit pure tensor
section and gives the same norm identity there.
\end{proof}
